\documentclass[11pt,letterpaper]{article}
\usepackage{nl_tps_common}
\hypersetup{
  pdftitle={GCPL / NL-TPS Defect Resolution Register},
  pdfsubject={Controlled dispositions for the 19 August 2026 corpus and tooling defect review}
}

\begin{document}
\NLSetPageStyle{GCPL / NL-TPS DEFECT RESOLUTION REGISTER}{Second-pass corpus and tooling review}
\fancyhead[R]{\small\bfseries\color{NLGOLD}CORRECTIVE ACTION}
\NLTitleBlock{CONTROLLED CORRECTIVE-ACTION RECORD}{Defect Resolution Register}{Governed Clinical Planning Language (GCPL) / NL-TPS}{\begin{minipage}{0.95\textwidth}
\NLMeta{Register identifier}{DRR-001}
\NLMeta{Source}{Second-pass defect review and RRR-001 repository-readiness reconciliation}
\NLMeta{Document version}{0.2 - Repository-readiness disposition baseline}
\NLMeta{Date}{19 August 2026}
\NLMeta{Coverage}{16 defects plus 12 repository-readiness recommendations}
\NLMeta{Status}{Corrective actions recorded; three controlled residuals remain open}
\end{minipage}}{This register records repository corrections and their structural verification. It does not approve a clinical value, accept residual risk, execute a V\&V check, commission an interface or model, authorize patient use, or replace named review by the accountable radiation-oncology professionals and quality authorities.}

\tableofcontents
\clearpage
\ifdefined\NLTPSCombinedDocument\else\pagenumbering{arabic}\fi

\NLFrontMatterSection{Document control}

\begin{small}
\begin{longtable}{@{}L{0.20\textwidth}L{0.72\textwidth}@{}}
\toprule
\rowcolor{NLTableHeader}\textbf{Field} & \textbf{Value} \\
\midrule
\endfirsthead
\toprule
\rowcolor{NLTableHeader}\textbf{Field} & \textbf{Value} \\
\midrule
\endhead
\midrule
\multicolumn{2}{r}{\scriptsize Continued on next page} \\
\endfoot
\bottomrule
\endlastfoot
Owner & GCPL / NL-TPS Program, with T-SYS corrective-action coordination and accountable clinical, physics, quality, security, interface, and independent V\&V review as applicable \\
\addlinespace[2pt]
Machine-readable source & \nolinkurl{spec/defects.yaml} \\
\addlinespace[2pt]
Controlled baseline & Requirements corpus v0.1; documentation suite v0.4; ENG-PKG-01 v0.6; ADR-001 and ADR-002 v0.2 \\
\addlinespace[2pt]
Change rule & A disposition changes only with reviewed source, evidence, regression results, residual-action owner, and impact analysis across requirements, hazards, interfaces, components, risks, V\&V, and release gates. \\
\addlinespace[2pt]
Closure meaning & ``Corrected'' means that the repository correction and its structural regression evidence exist. It never means that clinical risk is accepted, a V\&V claim passed, or clinical release is authorized. \\
\end{longtable}
\end{small}

\NLFrontMatterSection{Disposition summary}

\begin{nlnotice}{FAIL-CLOSED RELEASE POSITION}
Gate 0 structural checks may pass. Gate 1 and all later clinical or release gates remain blocked until the 119 HLR risk-score records and every applicable quantitative quality target have approved values and evidence. The active checkout also remains in OneDrive until a deliberate, verified migration is performed. No correction in this register grants clinical authority.
\end{nlnotice}

\begin{small}
\begin{longtable}{@{}L{0.07\textwidth}L{0.12\textwidth}L{0.31\textwidth}L{0.42\textwidth}@{}}
\toprule
\rowcolor{NLTableHeader}\textbf{ID} & \textbf{Severity} & \textbf{Disposition} & \textbf{Result} \\
\midrule
\endfirsthead
\toprule
\rowcolor{NLTableHeader}\textbf{ID} & \textbf{Severity} & \textbf{Disposition} & \textbf{Result} \\
\midrule
\endhead
\midrule
\multicolumn{4}{r}{\scriptsize Continued on next page} \\
\endfoot
\bottomrule
\endlastfoot
A1 & Critical & Corrected & All 2,144 controlled entities and V\&V claims carry one or more hazard links. \\
\addlinespace[2pt]
A2 & Critical & Corrected & All 2,088 risks derive failure, effect, and control from the hazard catalog; full normative text is retained. \\
\addlinespace[2pt]
A3 & Critical & Engineering control complete; clinical closure pending & 119 parent score records and inheritance rules exist; Gate 1 rejects unapproved scores. \\
\addlinespace[2pt]
A4 & Major & Engineering control complete; owner approval pending & Ten quantitative target classes exist; Gate 1 rejects missing approved values. \\
\addlinespace[2pt]
A5 & Major & Corrected & Twenty-five records are safety and assurance constraints; legacy aliases remain stable. \\
\addlinespace[2pt]
A6 & Major & Corrected & Twenty-nine ConOps seeds use the \texttt{CONOPS-} namespace and typed \texttt{SUPERSEDED\_BY} relations. \\
\addlinespace[2pt]
A7 & Major & Corrected & Applicable requirements allocate explicitly to registration, segmentation, dose, proton, intent, and architecture adapters. \\
\addlinespace[2pt]
B1 & Major & Corrected & Escaped ampersands no longer split controlled LaTeX rows; a negative unit test covers the case. \\
\addlinespace[2pt]
B2 & Major & Corrected & HLR bundle provenance is parsed from source metadata and bound to its hash. \\
\addlinespace[2pt]
B3 & Major & Corrected in source; hosted run pending push & Push and pull-request CI executes structural gates and regression tests. \\
\addlinespace[2pt]
B4 & Major & Corrected & Enforcement tools consume controlled architecture and terminology data. \\
\addlinespace[2pt]
B5 & Major & Corrected & HLR and traceability JSON Schemas are enforced without an additional local package. \\
\addlinespace[2pt]
B6 & Major & Corrected & Exact MPS dependency direction and global root/concept ownership are checked. \\
\addlinespace[2pt]
B7 & Major & Corrected & The planned Stage A specification artifacts, schemas, render guidance, and test guidance exist and reconcile. \\
\addlinespace[2pt]
C1 & Operational & BLOCKED\_ENVIRONMENT & Mandatory gate, not waived and not accepted. A strict location gate and managed Git worktree procedure exist; the synchronized checkout continues to fail by design until relocation. \\
\addlinespace[2pt]
C2 & Operational & Corrected & Generated outputs are ignored and previously tracked build products are removed from the source index. \\
\addlinespace[2pt]
K1 & Blocker & Corrected & Component records carry layer, name, realization responsibility, and lead team as typed fields instead of whichever column was longest. \\
\addlinespace[2pt]
K2 & Blocker & Corrected & Component, requirement, and sub-requirement construction schemas define and enforce what a constructed record contains. \\
\addlinespace[2pt]
K3 & Major & Engineering control complete; decision pending & Forty-three component overlap groups carry derived evidence and a proposed relation; approval is blocked until D-ENG-004 is decided. \\
\addlinespace[2pt]
K4 & Major & Corrected & D-ENG-001 through D-ENG-010 are a controlled register that names the gates and entities each decision blocks. \\
\addlinespace[2pt]
K5 & Major & Engineering control complete; review pending & Hazard specificity is explicit and a reviewed source-explicit set is a precondition for reviewed or approved construction state. \\
\addlinespace[2pt]
K6 & Moderate & Corrected & MQA realization requirements, sub-requirements, and subassemblies are typed apart and reconciled per key rather than by sum. \\
\addlinespace[2pt]
K7 & Major & Engineering control complete; approval pending & Each entity type carries a recorded hand, derived, or inherited construction treatment with a named-approver exception path. \\
\end{longtable}
\end{small}

\section{Construction-readiness findings}

\subsection{K1 and K2 --- Records that can carry a specification}

The component catalogs share one five-column layout of identifier, layer, name,
realization responsibility, and lead team. The trace importer previously kept whichever
column was longest, which resolved to the layer label for all 78 category components and
to the name for the 45 core and 61 interface components, so no component record carried
the responsibility that construction elaborates. The importer now selects the catalog row
explicitly and captures four typed fields; the eight machine-QA subassemblies are handled
the same way, including their recorded component allocations. The correction propagates
into the risk register, whose component records previously carried a layer label as their
requirement context.

Construction schemas now define what a constructed component, requirement, or
sub-requirement contains, with unknown fields rejected and conditional rules requiring an
interface contract, failure and degraded behavior, an acceptance criterion, and a named
reviewer before a record may claim reviewed or approved state. One exemplar of each class
is committed so the field set is reviewed at three records rather than at four hundred.

\subsection{K3, K4, and K7 --- Decisions that constrain construction}

Component overlap across the core, interface, and category views is derived rather than
asserted: 43 groups, of which two share a responsibility verbatim, seven share a name, and
thirty-four share only an identifier stem. Each group carries a proposed relation and stays
pending review, and no component in an undecided group may be approved. The ten open
engineering decisions are now machine-readable with explicit blocked gates and entities, so
the plan's work-permitted-meanwhile column is enforced rather than intended. The
construction treatment policy assigns hand, derived, or inherited treatment to each entity
type, retaining every identifier and every verification claim while directing individual
review to the records that carry engineering judgement.

\subsection{K5 and K6 --- Discrimination and taxonomy}

Interface and sub-interface requirements averaged 13.8 of 18 hazards, all machine-QA records
shared one set, and 3 percent of the graph carried a source-explicit mapping. Hazard
refinement is therefore made a construction obligation rather than a separate project: a
record cannot reach reviewed or approved state on an inherited hazard set. Machine-QA
realization requirements, sub-requirements, and subassemblies are typed apart, and entity
counts reconcile per key rather than by sum, which is the check the previous taxonomy
defect passed.

\section{Critical findings}

\subsection{A1 --- Complete hazard trace}

\textbf{Correction.} The eighteen ConOps hazards are controlled in \nolinkurl{spec/hazards.yaml}; allocation policy is controlled in \nolinkurl{spec/allocations.yaml}. The deterministic graph in \nolinkurl{mps/import/traceability.json} contains all 2,144 HLR, sub-requirement, HLIR, SIR, core-component, interface-component, category-component, and Machine-QA entities. Each record has one or more hazard IDs and a unique V\&V claim ID. The V\&V generator prints those hazards on every check.

\textbf{Verification.} \nolinkurl{tools/spec/build_trace_graph.py --check} shall report exactly 2,144 hazard-linked entities and claims. The aggregate spec gate schema-validates the graph. This is structural trace evidence, not a clinical judgment that the allocation is sufficient.

\subsection{A2 --- Hazard-derived risk semantics and lossless source text}

\textbf{Correction.} The risk builder obtains failure conditions, potential effects, and control strategies from the linked hazard rather than an interface-family sentence pool. Full normative text and its SHA-256 are retained in \nolinkurl{spec/risks.yaml}. The PDF table may use a visibly marked \texttt{\string\ldots\{\}} excerpt, but the structured record is not truncated. GOV-001 is linked to H-13 and H-18 and therefore addresses authority representation and automation bias rather than DICOM conformance.

\textbf{Verification.} \nolinkurl{tools/spec/build_risk_register.py --check} shall reconcile 2,088 records, reject the data-losing four-dot pattern, and confirm HLR score inheritance.

\subsection{A3 --- Governed scoring without fabricated clinical values}

\textbf{Correction.} \nolinkurl{spec/risk_scores.yaml} contains exactly one score record for each of 119 HLRs. Children inherit an approved parent score unless an approved exception exists; components use the maximum approved linked-HLR score unless an approved override exists. The 2,088 records are not separately hand-scored.

\textbf{Residual action.} All S/O/D values, RPN calculations, mitigation-effectiveness findings, residual-risk decisions, and acceptability decisions remain pending a facilitated multidisciplinary review. Gate 1 fails closed until the 119 records are approved. This finding is not clinically closed.

\section{Major corpus findings}

\subsection{A4 and A5 --- Separate quality targets from safety and assurance constraints}

The 25 legacy HNFR/NFSR/NFC records are canonically classified as cross-cutting safety and assurance constraints while their identifiers remain valid trace aliases. Separately, \nolinkurl{spec/quality_attributes.yaml} defines target classes for response latency, durable-job time, capacity, availability, recovery time and point, retention, DICOM transactions, human factors, and model performance or drift. Values are deliberately null until intended use, site workflow, evidence, and named owners establish them. Gate 1 rejects an applicable target without an approved value, unit, scope, evidence source, owner, and acceptance method.

\subsection{A6 --- Unambiguous source-to-canonical relations}

The 29 ConOps seeds now use the \texttt{CONOPS-} namespace. \nolinkurl{spec/requirements.yaml} represents each mapping as \texttt{SUPERSEDED\_BY}; \nolinkurl{tools/spec/check_seed_namespace.py} checks uniqueness, relation type, target existence, and agreement with the rendered HLR crosswalk. No canonical HLR ID was renamed.

\subsection{A7 --- Clinical adapter allocation}

The core realization now allocates applicable child requirements beyond the generic planning orchestrator to the performing adapter. The enforced allocation report is C-ARCH-01: 19; C-INTENT-01: 22; C-REG-01: 6; C-SEG-01: 20; C-PBT-01: 6; and C-DOSE-01: 33. These are logical allocations, not claims of implementation, DICOM conformance, commissioning, or clinical validation.

\section{Major tooling findings}

\subsection{B1 and B2 --- Reliable parsing and provenance}

All controlled table parsers separate columns with \texttt{(?<!\\)\&}, so an escaped \texttt{Q\string\&A} remains within its normative field. Import metadata is parsed from \texttt{\string\NLMeta} fields, normalized, and bound to the source-file hash. Unit tests cover the escaped-ampersand, metadata, and invalid-schema cases.

\subsection{B3 --- Continuous gates}

\nolinkurl{.github/workflows/controlled-spec-gates.yml} runs on pushes and pull requests. It executes regression tests, location constraints, deterministic HLR import, MPS layering, namespace checks, reclassification, clinical allocation, hazard trace, hazard-derived risk, schema reconciliation, Gate 0 readiness, and V\&V generated-document drift. The workflow definition is present; hosted execution evidence can exist only after the change is pushed to GitHub.

\subsection{B4 and B5 --- Specification-as-data and schema enforcement}

The HLR builder and MPS checker consume \nolinkurl{spec/architecture.yaml} and \nolinkurl{spec/terminology.yaml}. The aggregate gate reconciles baseline identity, counts, hazard families, trace claims, risks, and defect records. The repository provides a dependency-free validator for the committed JSON Schema subset and applies it to both the HLR mirror and trace graph.

\subsection{B6 --- MPS language boundaries}

\nolinkurl{spec/architecture.yaml} controls the permitted dependency matrix, file-per-root persistence, disjoint root/non-root concept collections, global ownership scope, and prohibited data or credentials. The checker enforces exact dependency direction and global single ownership across four languages and 65 concepts. This does not represent the live MPS project or generators as complete.

\subsection{B7 --- Required specification artifacts}

The Stage A specification set now includes architecture, terminology, decisions, requirements, hazards, interfaces, components, allocations, V\&V, risks, risk scores, quality attributes, defect dispositions, schemas, render guidance, and test guidance. These files constrain and test implementation; under ADR-001 they remain mirrors and supplements until an approved Stage C authority cutover.

\section{Operational findings}

\subsection{C1 --- Workspace location}

C1 is dispositioned \texttt{BLOCKED\_ENVIRONMENT}. It is a known environmental prerequisite, not a model exception: it is not waived, not accepted, and no exception mechanism applies to it. The gate is mandatory and shall not be relaxed, skipped, or granted a permitted-failure allowance in any workflow.

The synchronized checkout is inside OneDrive and its path contains whitespace, so \nolinkurl{tools/repo/check_workspace_location.py --enforce} rejects it and will continue to do so by design. \nolinkurl{scripts/create_mps_worktree.ps1} requires a clean worktree, refuses an existing, synchronized, or whitespace-bearing destination, creates or attaches a dedicated materialization branch, verifies the shared Git common directory and source commit, and does not delete the source. Because the checker resolves the workspace root from the executing tree, the gate turns green in the managed worktree while remaining red in the synchronized checkout.

Acceptance is every gate green from the relocated worktree, not the prior count plus an explanation. One residual is explicit: the managed worktree relocates the working tree only. The controlling Git common directory remains in the synchronized location, so relocating the object store is a separate decision.

\subsection{C2 --- Generated-artifact control}

The ignore policy now covers \nolinkurl{output/}, \nolinkurl{tmp/pdfs/}, Python bytecode, and prior literal build locations. Generated PDFs and LaTeX intermediates are removed from the Git index but may remain locally for review. ART-POL-001 requires CI or an explicit release process, stable names, source hashes, toolchain identity, output hashes, page counts, release status, and a provenance manifest.

\section{Repository-readiness review reconciliation}

RRR-001 reconciles the twelve follow-on recommendations made against GitHub commit \texttt{399244f} with the newer local corrective baseline. It records a repository correction as complete only within its structural scope and keeps external or approval-dependent work explicit.

\begin{small}
\begin{longtable}{@{}L{0.10\textwidth}L{0.12\textwidth}L{0.34\textwidth}L{0.36\textwidth}@{}}
\toprule
\rowcolor{NLTableHeader}\textbf{ID} & \textbf{Priority} & \textbf{Disposition} & \textbf{Residual action} \\
\midrule
\endfirsthead
\toprule
\rowcolor{NLTableHeader}\textbf{ID} & \textbf{Priority} & \textbf{Disposition} & \textbf{Residual action} \\
\midrule
\endhead
\midrule
\multicolumn{4}{r}{\scriptsize Continued on next page} \\
\endfoot
\bottomrule
\endlastfoot
RRR-001 & P0 & Tracked build products removed; hygiene scanning added & Historical path disclosure remains until an owner-approved history decision. \\
\addlinespace[2pt]
RRR-002 & P0 & Current suite review artifact and provenance manifest generated & Review artifact is not a governed or clinical release. \\
\addlinespace[2pt]
RRR-003 & P0 & Structural, hygiene, schema, drift, standalone-document, and suite CI defined & First hosted execution evidence follows commit and push. \\
\addlinespace[2pt]
RRR-004 & P0 & Managed space-free Git worktree procedure defined & Run only after the controlling worktree is clean. \\
\addlinespace[2pt]
RRR-005 & P1 & Stage A materialization handoff and evidence checklist defined & Live MPS project, languages, importer, generators, editors, and tests remain not started. \\
\addlinespace[2pt]
RRR-006 & P1 & Root and non-root concept collections made disjoint and schema-enforced & Confirm the same invariant in live MPS. \\
\addlinespace[2pt]
RRR-007 & P1 & Stable schema URNs and architecture, decision, terminology, skeleton, HLR, and trace validation added & Stage A validation does not authorize Stage C. \\
\addlinespace[2pt]
RRR-008 & P1 & Explicit approval state and fail-closed Stage C/release checker added & Named approval arrays remain empty. \\
\addlinespace[2pt]
RRR-009 & P1 & ART-POL-001 selects CI/release artifacts with a provenance manifest & Governed release authorization remains false. \\
\addlinespace[2pt]
RRR-010 & P1/P2 & Root governance documents added; Word binaries quarantined & Repository owner and institutional license decision remains pending. \\
\addlinespace[2pt]
RRR-011 & P2 & V\&V generator \texttt{-Check} drift mode present & V\&V claims remain planned and unexecuted. \\
\addlinespace[2pt]
RRR-012 & P2 & Corpus, ADR supplement, and suite assembly versions stated separately & None within version-labeling scope. \\
\end{longtable}
\end{small}

\section{Verification commands and release interpretation}

\begin{small}
\begin{longtable}{@{}L{0.47\textwidth}L{0.45\textwidth}@{}}
\toprule
\rowcolor{NLTableHeader}\textbf{Command} & \textbf{Acceptance interpretation} \\
\midrule
\endfirsthead
\toprule
\rowcolor{NLTableHeader}\textbf{Command} & \textbf{Acceptance interpretation} \\
\midrule
\endhead
\midrule
\multicolumn{2}{r}{\scriptsize Continued on next page} \\
\endfoot
\bottomrule
\endlastfoot
{\raggedright\ttfamily python -m unittest discover -s tests -v\par} & Tooling regressions pass; this is not clinical V\&V. \\
\addlinespace[2pt]
{\raggedright\ttfamily python tools/mps/build\_hlr\_import\_bundle.py --check\par} & The 119-HLR Stage A mirror is schema-valid and deterministic. \\
\addlinespace[2pt]
{\raggedright\ttfamily python tools/spec/build\_trace\_graph.py --check\par} & All 2,144 entities and claim IDs have a materialized hazard allocation. \\
\addlinespace[2pt]
{\raggedright\ttfamily python tools/spec/build\_risk\_register.py --check --gate 0\par} & The 2,088 risk records are hazard-derived and structurally consistent. \\
\addlinespace[2pt]
{\raggedright\ttfamily python tools/spec/check\_controlled\_specs.py\par} & Required specs, counts, schemas, and baselines reconcile. \\
\addlinespace[2pt]
{\raggedright\ttfamily python tools/spec/check\_gate\_readiness.py --gate 0\par} & Nonclinical structural work may proceed within the frozen architecture. \\
\addlinespace[2pt]
{\raggedright\ttfamily python tools/spec/check\_gate\_readiness.py --gate 1\par} & Expected to fail until named owners approve all required risk scores and quality targets. \\
\end{longtable}
\end{small}

\section{Controlled residual-action list}

\begin{enumerate}
\item \textbf{A3 owner action:} convene the multidisciplinary hazard/FMEA review, approve the scoring scale, complete 119 HLR scores, document exceptions, and record residual-risk decisions.
\item \textbf{A4 owner action:} approve applicable quantitative values, units, scopes, evidence sources, and acceptance methods for the selected intended use and implementation slice.
\item \textbf{C1 user action:} after the corrective work is committed or stashed and reviewed, run the managed-worktree script for a space-free, non-synchronized destination; verify the registered worktree before opening it in MPS.
\item \textbf{B3 evidence action:} after push, retain the first successful hosted CI run as repository evidence; investigate any platform-specific difference rather than waiving it silently.
\end{enumerate}

\begin{nlnotice}{NO CLINICAL CLOSURE BY DOCUMENTATION}
This register may close a software or documentation defect only within its stated evidence scope. It cannot establish clinical efficacy, safety, risk acceptability, independent V\&V acceptance, DICOM conformance, commissioning, accreditation, regulatory status, or permission to use GCPL / NL-TPS for a patient.
\end{nlnotice}

\end{document}
